% title:          The inekoalaty game
% area:           games, algebraic-inequalities
% methods:
% difficulty:
% difficulty-ai:  hard
% status:
% review:         none
% --- history: all optional, leave EMPTY if unknown ---
% origin:         imo
% origin-date:    2025-07-16
% origin-number:  5
% also-in:
% taken-from:
% references:     IMO 2025 (Sunshine Coast, Australia), Day 2 (16 July 2025), Problem 5, proposed by Italy - https://www.imo-official.org/problems/2025/ ; official English paper - https://www.imo-official.org/assets/documents/problems/2025/2025_eng.pdf ; IMO Shortlist 2025, A3 (Italy), with players named Coral and Joey, with official solution (answer: Alice wins iff lambda > 1/sqrt2, Bazza wins iff lambda < 1/sqrt2, neither at lambda = 1/sqrt2) - https://www.imo-official.org/assets/documents/problems/2025/IMO2025SL.pdf ; AoPS Wiki, 2025 IMO Problems/Problem 5 (archived) - https://web.archive.org/web/20250925165721/https://artofproblemsolving.com/wiki/index.php/2025_IMO_Problems/Problem_5
% history:        ai
\begin{problem}
Alice and Bazza are playing the \textit{inekcalaty} game, a two-player game whose rules depend on a positive real number \(\lambda\) which is known to both players. On the \(n\)th turn of the game (starting with \(n = 1\)) the following happens:

\begin{itemize}
    \item If \(n\) is odd, Alice chooses a nonnegative real number \(x_n\) such that \(x_1 + x_2 + \cdots + x_n \leq \lambda n\).
    
    \item If \(n\) is even, Bazza chooses a nonnegative real number \(x_n\) such that \(x_1^2 + x_2^2 + \cdots + x_n^2 \leq n\).
\end{itemize}

If a player cannot choose a suitable \(x_m\) the game ends and the other player wins. If the game goes on forever, neither player wins. All chosen numbers are known to both players.

Determine all values of \(\lambda\) for which Alice has a winning strategy and all those for which Bazza has a winning strategy.
\end{problem}
